\section{The category of affine schemes}
\label{section-category-affine-schemes}

\noindent
Note that if $Y$ is an affine scheme, then its
points are in canonical $1-1$ bijection with
prime ideals in $\Gamma(Y, \mathcal{O}_Y)$.

\begin{lemma}
\label{lemma-morphism-into-affine-where-point-goes}
Let $X$ be a locally ringed space.
Let $Y$ be an affine scheme.
Let $f \in \Mor(X, Y)$ be a morphism
of locally ringed spaces. Given a point $x \in X$
consider the ring maps
$$
\Gamma(Y, \mathcal{O}_Y) \xrightarrow{f^\sharp}
\Gamma(X, \mathcal{O}_X) \to \mathcal{O}_{X, x}
$$
Let $\mathfrak p \subset \Gamma(Y, \mathcal{O}_Y)$ denote
the inverse image of $\mathfrak m_x$. Let $y \in Y$ be the
corresponding point. Then $f(x) = y$.
\end{lemma}

\begin{proof}
Consider the commutative diagram
$$
\xymatrix{
\Gamma(X, \mathcal{O}_X) \ar[r] &
\mathcal{O}_{X, x} \\
\Gamma(Y, \mathcal{O}_Y) \ar[r] \ar[u] &
\mathcal{O}_{Y, f(x)} \ar[u]
}
$$
(see the discussion of $f$-maps below
Sheaves, Definition \ref{sheaves-definition-f-map}).
Since the right vertical arrow is local
we see that $\mathfrak m_{f(x)}$ is the
inverse image of $\mathfrak m_x$. The result
follows.
\end{proof}

\begin{lemma}
\label{lemma-f-open}
Let $X$ be a locally ringed space.
Let $f \in \Gamma(X, \mathcal{O}_X)$.
The set
$$
D(f) = \{x \in X \mid \text{image }f \not\in \mathfrak m_x\}
$$
is open. Moreover $f|_{D(f)}$ has an inverse.
\end{lemma}

\begin{proof}
This is a special case of Modules, Lemma \ref{modules-lemma-s-open}, but
we also give a direct proof.
Suppose that $U \subset X$ and $V \subset X$ are
two open subsets such that $f|_U$ has an inverse
$g$ and $f|_V$ has an inverse $h$. Then clearly
$g|_{U\cap V} = h|_{U\cap V}$. Thus it suffices
to show that $f$ is invertible in an open neighbourhood
of any $x \in D(f)$. This is clear because
$f \not \in \mathfrak m_x$ implies that $f \in \mathcal{O}_{X, x}$
has an inverse $g \in \mathcal{O}_{X, x}$ which means there
is some open neighbourhood $x \in U \subset X$ so
that $g \in \mathcal{O}_X(U)$ and $g\cdot f|_U = 1$.
\end{proof}

\begin{lemma}
\label{lemma-f-open-affine}
In Lemma \ref{lemma-f-open} above, if $X$ is an affine scheme,
then the open $D(f)$ agrees with the standard open $D(f)$
defined previously (in
Algebra, Definition \ref{algebra-definition-spectrum-ring}).
\end{lemma}

\begin{proof}
Omitted.
\end{proof}

\begin{lemma}
\label{lemma-morphism-into-affine}
\begin{reference}
A reference for this fact is \cite[II, Err 1, Prop. 1.8.1]{EGA}
where it is attributed to J. Tate.
\end{reference}
Let $X$ be a locally ringed space.
Let $Y$ be an affine scheme.
The map
$$
\Mor(X, Y)
\longrightarrow
\Hom(\Gamma(Y, \mathcal{O}_Y), \Gamma(X, \mathcal{O}_X))
$$
which maps $f$ to $f^\sharp$ (on global sections) is bijective.
\end{lemma}

\begin{proof}
Since $Y$ is affine we have
$(Y, \mathcal{O}_Y) \cong (\Spec(R), \mathcal{O}_{\Spec(R)})$
for some ring $R$.
During the proof we will use facts about $Y$ and
its structure sheaf which are direct consequences of things
we know about the spectrum of a ring, see e.g.\ Lemma
\ref{lemma-spec-sheaves}.

\medskip\noindent
Motivated by the lemmas above we construct the inverse map.
Let $\psi_Y : \Gamma(Y, \mathcal{O}_Y) \to \Gamma(X, \mathcal{O}_X)$
be a ring map. First, we define the corresponding map of
spaces
$$
\Psi : X \longrightarrow Y
$$
by the rule of
Lemma \ref{lemma-morphism-into-affine-where-point-goes}.
In other words, given $x \in X$ we define $\Psi(x)$
to be the point of $Y$ corresponding to the prime
in $\Gamma(Y, \mathcal{O}_Y)$ which is the inverse
image of $\mathfrak m_x$ under the composition
$
\Gamma(Y, \mathcal{O}_Y) \xrightarrow{\psi_Y}
\Gamma(X, \mathcal{O}_X) \to
\mathcal{O}_{X, x}
$.

\medskip\noindent
We claim that the map $\Psi : X \to Y$ is continuous.
The standard opens $D(g)$, for $g \in \Gamma(Y, \mathcal{O}_Y)$
are a basis for the topology of $Y$. Thus it suffices to prove
that $\Psi^{-1}(D(g))$ is open. By construction of $\Psi$
the inverse image $\Psi^{-1}(D(g))$ is exactly the set
$D(\psi_Y(g)) \subset X$ which is open by Lemma \ref{lemma-f-open}.
Hence $\Psi$ is continuous.

\medskip\noindent
Next we construct a $\Psi$-map of sheaves from
$\mathcal{O}_Y$ to $\mathcal{O}_X$. By
Sheaves, Lemma \ref{sheaves-lemma-f-map-basis-below-structures}
it suffices to define ring maps
$\psi_{D(g)} : \Gamma(D(g), \mathcal{O}_Y) \to
\Gamma(\Psi^{-1}(D(g)), \mathcal{O}_X)$
compatible with restriction maps.
We have a canonical isomorphism
$\Gamma(D(g), \mathcal{O}_Y) = \Gamma(Y, \mathcal{O}_Y)_g$,
because $Y$ is an affine scheme.
Because $\psi_Y(g)$ is invertible on $D(\psi_Y(g))$
we see that there is a canonical map
$$
\Gamma(Y, \mathcal{O}_Y)_g
\longrightarrow
\Gamma(\Psi^{-1}(D(g)), \mathcal{O}_X)
=
\Gamma(D(\psi_Y(g)), \mathcal{O}_X)
$$
extending the map $\psi_Y$
by the universal property of localization.
Note that there is no choice but to take the canonical map here!
And we take this, combined
with the canonical identification
$\Gamma(D(g), \mathcal{O}_Y) = \Gamma(Y, \mathcal{O}_Y)_g$, to
be $\psi_{D(g)}$. This is compatible with localization since the
restriction mapping on the affine schemes are defined in terms
of the universal properties of localization also, see
Lemmas \ref{lemma-spec-sheaves} and \ref{lemma-standard-open}.

\medskip\noindent
Thus we have defined a morphism of ringed spaces
$(\Psi, \psi) : (X, \mathcal{O}_X) \to (Y, \mathcal{O}_Y)$
recovering $\psi_Y$ on global sections. To see that it is
a morphism of locally ringed spaces we have to show that
the induced maps on local rings
$$
\psi_x : \mathcal{O}_{Y, \Psi(x)} \longrightarrow \mathcal{O}_{X, x}
$$
are local. This follows immediately from the commutative diagram
of the proof of Lemma \ref{lemma-morphism-into-affine-where-point-goes}
and the definition of $\Psi$.

\medskip\noindent
Finally, we have to show that the constructions
$(\Psi, \psi) \mapsto \psi_Y$ and the construction
$\psi_Y \mapsto (\Psi, \psi)$ are inverse to each other.
Clearly, $\psi_Y \mapsto (\Psi, \psi) \mapsto \psi_Y$.
Hence the only thing to prove is that given $\psi_Y$
there is at most one pair $(\Psi, \psi)$ giving rise
to it. The uniqueness of $\Psi$ was shown in Lemma
\ref{lemma-morphism-into-affine-where-point-goes} and
given the uniqueness of $\Psi$ the uniqueness of the
map $\psi$ was pointed out during the course of the proof
above.
\end{proof}

\begin{lemma}
\label{lemma-category-affine-schemes}
The category of affine schemes is equivalent to the opposite of the
category of rings. The equivalence is given by the functor that associates
to an affine scheme the global sections of its structure sheaf.
\end{lemma}

\begin{proof}
This is now clear from Definition \ref{definition-affine-scheme}
and Lemma \ref{lemma-morphism-into-affine}.
\end{proof}

\begin{lemma}
\label{lemma-standard-open-affine}
Let $Y$ be an affine scheme.
Let $f \in \Gamma(Y, \mathcal{O}_Y)$.
The open subspace $D(f)$ is an affine scheme.
\end{lemma}

\begin{proof}
We may assume that $Y = \Spec(R)$ and $f \in R$.
Consider the morphism of affine schemes
$\phi : U = \Spec(R_f) \to \Spec(R) = Y$ induced by the ring
map $R \to R_f$. By Algebra, Lemma \ref{algebra-lemma-standard-open}
we know that it is a homeomorphism onto $D(f)$.
On the other hand, the map $\phi^{-1}\mathcal{O}_Y \to \mathcal{O}_U$
is an isomorphism on stalks, hence an isomorphism. Thus we see
that $\phi$ is an open immersion. We conclude that $D(f)$
is isomorphic to $U$ by Lemma \ref{lemma-open-immersion}.
\end{proof}

\begin{lemma}
\label{lemma-fibre-product-affine-schemes}
The category of affine schemes has finite products, and fibre products.
In other words, it has finite limits. Moreover, the products
and fibre products in the category of affine schemes
are the same as in the category of locally ringed spaces.
In a formula, we have (in the category of locally ringed spaces)
$$
\Spec(R) \times \Spec(S) =
\Spec(R \otimes_{\mathbf{Z}} S)
$$
and given ring maps $R \to A$, $R \to B$ we have
$$
\Spec(A) \times_{\Spec(R)} \Spec(B)
=
\Spec(A \otimes_R B).
$$
\end{lemma}

\begin{proof}
This is just an application of Lemma \ref{lemma-morphism-into-affine}.
First of all, by that lemma, the affine scheme
$\Spec(\mathbf{Z})$ is the final object in the category
of locally ringed spaces. Thus the first displayed formula
follows from the second. To prove the second note that
for any locally ringed space $X$ we have
\begin{eqnarray*}
\Mor(X, \Spec(A \otimes_R B))
& = &
\Hom(A \otimes_R B, \mathcal{O}_X(X)) \\
& = &
\Hom(A, \mathcal{O}_X(X))
\times_{\Hom(R, \mathcal{O}_X(X))}
\Hom(B, \mathcal{O}_X(X)) \\
& = &
\Mor(X, \Spec(A))
\times_{\Mor(X, \Spec(R))}
\Mor(X, \Spec(B))
\end{eqnarray*}
which proves the formula.
See Categories, Section \ref{categories-section-fibre-products} for the
relevant definitions.
\end{proof}

\begin{lemma}
\label{lemma-disjoint-union-affines}
Let $X$ be a locally ringed space.
Assume $X = U \amalg V$ with $U$ and $V$ open and
such that $U$, $V$ are affine schemes. Then $X$ is an affine scheme.
\end{lemma}

\begin{proof}
Set $R = \Gamma(X, \mathcal{O}_X)$.
Note that $R = \mathcal{O}_X(U) \times \mathcal{O}_X(V)$
by the sheaf property. By Lemma \ref{lemma-morphism-into-affine}
there is a canonical morphism of locally ringed spaces
$X \to \Spec(R)$. By Algebra, Lemma \ref{algebra-lemma-spec-product}
we see that as a topological space
$\Spec(\mathcal{O}_X(U)) \amalg \Spec(\mathcal{O}_X(V)) =
\Spec(R)$
with the maps coming from the ring homomorphisms
$R \to \mathcal{O}_X(U)$ and $R \to \mathcal{O}_X(V)$.
This of course means that $\Spec(R)$ is the coproduct
in the category of locally ringed spaces as well.
By assumption the morphism $X \to \Spec(R)$ induces an isomorphism
of $\Spec(\mathcal{O}_X(U))$ with $U$ and similarly
for $V$. Hence $X \to \Spec(R)$ is an isomorphism.
\end{proof}
